For board size $n=2L\ge4$, let $U$ count unlike-type orthogonal neighbour pairs and $V$ unlike-type diagonal neighbour pairs. The one-way attack count is $U+V$. Each adjacent $2\times2$ block has at most four unlike pairs among its six pairs; summing gives $2U+V\le4n^2$. Also $V\le2n^2$. Hence $U+V\le3n^2$, attained by alternating columns of rooks and bishops: every diagonal pair and half of the orthogonal pairs is unlike. Thus $M=12L^2$.
For the pentagon, let $b=ZO$, $a=HA$, and let $h,k$ be the distances of $D,ZO$ from line $HA$. Regular-pentagon similarity gives $\dfrac{h}{k}=1+\dfrac{a}{b}$ (both equal the golden ratio). Thus $bh=k(a+b)=2[ZHAO]=2M$. Adding the triangles above and below $ZO$ gives $[DZHAO]+[ORZ]=M+bh=3M$, so $T=36L^2$.
Since $301=7\cdot43$, the order of 3 is the least common multiple of its orders 6 and 42, namely 42. For 43, $3^{21}\equiv-1$, $3^{14}\equiv36$, $3^6\equiv41$, ruling out every proper divisor of 42. Hence $L=\lfloor \dfrac{T}{252}\rfloor=\lfloor \dfrac{L^2}{7}\rfloor$. This means $7L\le L^2<7L+7$, whose sole positive integer solution is 7. The exceptional two-by-two board at $L=1$ cannot work: it has fewer than 16 possible ordered distinct pairs, so $T<48$ and the required count is zero. Finally $\boxed{L=7,M=588,T=1764}$.
\begin{asy}
import graph;import geometry;usepackage("amsmath");size(7cm);
for(int r=0;r<4;++r)for(int c=0;c<4;++c){draw((c,r)--(c+1,r)--(c+1,r+1)--(c,r+1)--cycle);label(c%2==0?"R":"B",(c+.5,r+.5));}label("alternating columns",(2,-.5));pair D=(8,1),Z=(8-sin(pi*72/180),cos(pi*72/180)),H=(8-sin(pi*36/180),-cos(pi*36/180)),A=(8+sin(pi*36/180),-cos(pi*36/180)),O=(8+sin(pi*72/180),cos(pi*72/180)),R=(8,-1-2cos(pi*36/180));draw(D--Z--H--A--O--cycle);draw(O--R--Z,blue);draw(H--A,dashed);label("$D$",D,(0,1));label("$Z$",Z,(-1,0));label("$H$",H,(-1,0));label("$A$",A,(1,0));label("$O$",O,(1,0));label("$R$",R,(0,-1));
\end{asy}

